\section{Training Procedure and Optimization}
\label{sec:training}

The training phase of our classification architecture was conducted through an iterative process, continuously adapting both dataset composition and hyperparameter configurations to optimize model accuracy while working under strict hardware constraints.

\subsection{Dataset Curation and Iterative Expansion}
Initial training runs conducted on a small subset of data revealed a severe performance discrepancy compared to expected detector metrics. This was primarily attributed to insufficient sample density per class, as our initial subset contained only approximately 20 animals per breed. Although external resources such as the Cat API \cite{thecatapi} provide extensive image repositories, poor label quality resulted in an effective yield of only about 12 usable images per breed. To avoid manual labeling overhead, we expanded our data pipeline using established benchmark datasets.

First, we integrated the Oxford-IIIT Pet Dataset \cite{parkhi2012cats}, which provided an additional 2,744 images of target breeds. Visual inspection during integration revealed significant ambiguity between the \textit{Tabby} and \textit{Tiger Cat} categories. Recognizing that \textit{Tabby} represents a coat pattern rather than a distinct breed, we consolidated these samples under the \textit{Tiger Cat} label. To compensate for missing target breeds—specifically \textit{Golden Retriever}, \textit{German Shepherd}, \textit{Siberian Husky}, \textit{Dalmatian}, and \textit{Rottweiler}—we incorporated the Stanford Dogs Dataset \cite{khosla2011stanford}, securing approximately 150–200 images per breed. 

To further stabilize convergence on challenging classes like \textit{Tiger Cat}, we gathered 30 additional verified images per breed from open platforms (Unsplash.com and Pixabay.com) and progressively raised class density toward 50 images per breed. This targeted augmentation pushed validation accuracy beyond our initial baseline goal of 80\%. Finally, auditing unexpected validation drops identified a complete absence of \textit{Dalmatian} samples in early data folds; supplementing 359 validated \textit{Dalmatian} images resulted in a performance surge from 17\% to 69\% accuracy within 10 epochs.

\subsection{Hardware Constraints and Hyperparameter Tuning}
Initial training attempts utilizing an input resolution of $384 \times 384$ pixels with a batch size of 32 exceeded the memory limits of our 7.6\,GiB GPU. Memory allocation in Convolutional Neural Networks scales quadratically with spatial image dimensions: increasing input dimensions from $224 \times 224$ (50,176 pixels) to $384 \times 384$ (147,456 pixels) requires approximately three times the VRAM per image to store intermediate activation maps for backpropagation. 

To resolve this bottleneck, we constrained the spatial resolution to $224 \times 224$ pixels and reduced the batch size to 16. Halving the batch size linearly reduced peak VRAM consumption during the forward-backward pass, allowing stable parallel processing without GPU out-of-memory (OOM) faults.

\subsection{Mitigating Class Imbalance via Loss Weighting}
Given the residual variance in sample counts across breeds, unweighted training risks biasing the model toward over-represented majority classes to artificially minimize overall loss. To penalize misclassifications on minority classes without computationally expensive resampling, we integrated class-weighted cross-entropy loss using PyTorch's \texttt{nn.CrossEntropyLoss}.

For each class $i \in \{0, \dots, C-1\}$ (where $C=10$), weights $w_i$ are dynamically computed prior to training based on sample counts $n_i$ and total training samples $N_{\text{total}}$:
\begin{equation}
w_i = \frac{N_{\text{total}}}{C \cdot n_i}
\end{equation}
During backpropagation, errors on rare classes generate proportionally larger loss gradients (e.g., $w_i = 5.0$) compared to frequent classes (e.g., $w_i = 0.2$), strictly penalizing misclassifications on minority breeds and forcing balanced feature representation across all target outputs.

\subsection{Optimization, Convergence Dynamics, and Overfitting Control}
Optimization dynamics highlighted critical dependencies on learning rate selection. Excessively low learning rates caused the optimizer to stall on loss surface plateaus or get trapped in sub-optimal local minima due to insufficient step energy. Conversely, un-decayed learning rates induced erratic jumping around optimal minima during final epochs.

To prevent overfitting—where training accuracy continues to rise while validation performance degrades—we employed a multi-faceted regularization and monitoring strategy grounded in recent studies on generalization gap control \cite{zhang2024overfitguard}:
\begin{itemize}
    \item \textbf{Regularization \& Augmentation:} We applied Dropout layers alongside heavy data augmentation (Random Horizontal Flip, Random Rotation, Brightness/Contrast adjustments, and CutMix/CropMix transformations) to ensure the network rarely encountered identical pixel configurations twice.
    \item \textbf{Model Checkpointing \& Early Stopping:} To counter performance degradation over extended epoch runs, we logged validation accuracy at each epoch, saving model weights only when validation accuracy improved. Training was configured to terminate automatically if validation performance failed to improve over 10 consecutive epochs.
    \item \textbf{Generalization Gap Thresholding:} We actively monitored the divergence between training and validation accuracy ($\Delta_{\text{acc}} = \text{Acc}_{\text{train}} - \text{Acc}_{\text{val}}$) per epoch. Training was halted immediately if this accuracy gap exceeded a strict $7\%$ threshold ($\Delta_{\text{acc}} > 7\%$) or if epoch-over-epoch accuracy gains fell below $+0.1\%$, preventing the model from memorizing training-specific noise \cite{zhang2024overfitguard}.
\end{itemize}